% GENERATED FILE. Edit canonical lessons, then run npm run generate:books.
% canonical-source-hash: fnv1a64:a20581c5ca381d21
% canonical-lessons: UR-C534-bisat, UR-C534-mohra, UR-C534-piyada, UR-C534-mat, UR-C534-ludo

\chapter{On the Board}
\label{ch:ur-on-the-board}
\hlchaptermodality{534}

\begin{chapteropening}
\textbf{By the end of this chapter:} \emph{I can talk about board games and their pieces.}

Complete the last lesson of chapter 534: I can talk about board games and their pieces.
\end{chapteropening}

\section[\texorpdfstring{\ur{بساط} (bisāt)}{بساط (bisāt)}]{\ur{بساط} (bisāt) — a game board}

\label{lesson:UR-C534-bisat}



\begin{quote}
Before the new one: say the Urdu for a move (in a game), then the Urdu for to turn over.
\end{quote}

\subsection*{You'll want to know: \ur{بساط}}
\textbf{\ur{بساط}} — \emph{bisāt} — \textquotedblleft{}a game board\textquotedblright{}.

\begin{quote}
\ur{اگر} \ur{بساط} \ur{ہے} \ur{تو} \ur{ہم} \ur{شطرنج} \ur{کھیلیں} \ur{گے۔}
\end{quote}

\begin{quote}
\emph{agar bisāt hai to ham shatranj kheleṅ ge.}
\end{quote}

\begin{quote}
If you have a board, we'll play chess.
\end{quote}

Feminine.

\begin{cousinweb}
\textbf{\ur{بساط}} (\emph{bisāt}) is usually traced, through Persian, to Arabic \emph{bisāṭ}, \textquotedblleft{}a carpet, something spread out\textquotedblright{}.
\end{cousinweb}

\begin{grammarlens}[title={What you've built}]
One more word for saying what happens if.
\end{grammarlens}

\subsection*{Your turn}
\begin{itemize}
\raggedright
  \item \emph{Say it:} \emph{bisāt}
  \item \emph{Say it:} \emph{bisāt}, once more
  \item \emph{Say it:} \emph{agar bisāt hai to ham shatranj kheleṅ ge.} — If you have a board, we'll play chess.
  \item \emph{Recall:} say the Urdu for a move (in a game), then the Urdu for to turn over, then say \emph{bisāt} again
\end{itemize}

\begin{tcolorbox}[breakable,colback=teal!4,colframe=teal!35!black,title={Before you move on}]
What does \ur{بساط} mean? (\textquotedblleft{}A game board\textquotedblright{}.) Say it once more.
\end{tcolorbox}
\section[\texorpdfstring{\ur{مہرہ} (mohrā)}{مہرہ (mohrā)}]{\ur{مہرہ} (mohrā) — a piece (in a board game)}

\label{lesson:UR-C534-mohra}



\begin{quote}
Before the new one: say the Urdu for to turn over, then the Urdu for a game board.
\end{quote}

\subsection*{You'll want to know: \ur{مہرہ}}
\textbf{\ur{مہرہ}} — \emph{mohrā} — \textquotedblleft{}a piece (in a board game)\textquotedblright{}.

\begin{quote}
\ur{اگر} \ur{ایک} \ur{مہرہ} \ur{کھو} \ur{گیا} \ur{تو} \ur{کھیل} \ur{رک} \ur{جائے} \ur{گا۔}
\end{quote}

\begin{quote}
\emph{agar ek mohrā kho gayā to khel ruk jāe gā.}
\end{quote}

\begin{quote}
If a piece goes missing, the game will stop.
\end{quote}

Masculine.

\begin{cousinweb}
\textbf{\ur{مہرہ}} (\emph{mohrā}) is usually traced to Persian \emph{muhra}, \textquotedblleft{}a bead; a piece in a game\textquotedblright{}.
\end{cousinweb}

\begin{grammarlens}[title={What you've built}]
One more word for saying what happens if.
\end{grammarlens}

\subsection*{Your turn}
\begin{itemize}
\raggedright
  \item \emph{Say it:} \emph{mohrā}
  \item \emph{Say it:} \emph{mohrā}, once more
  \item \emph{Say it:} \emph{agar ek mohrā kho gayā to khel ruk jāe gā.} — If a piece goes missing, the game will stop.
  \item \emph{Recall:} say the Urdu for to turn over, then the Urdu for a game board, then say \emph{mohrā} again
\end{itemize}

\begin{tcolorbox}[breakable,colback=teal!4,colframe=teal!35!black,title={Before you move on}]
What does \ur{مہرہ} mean? (\textquotedblleft{}A piece (in a board game)\textquotedblright{}.) Say it once more.
\end{tcolorbox}
\section[\texorpdfstring{\ur{پیادہ} (piyādā)}{پیادہ (piyādā)}]{\ur{پیادہ} (piyādā) — a pawn (in chess)}

\label{lesson:UR-C534-piyada}



\begin{quote}
Before the new one: say the Urdu for a game board, then the Urdu for a piece (in a board game).
\end{quote}

\subsection*{You'll want to know: \ur{پیادہ}}
\textbf{\ur{پیادہ}} — \emph{piyādā} — \textquotedblleft{}a pawn (in chess)\textquotedblright{}.

\begin{quote}
\ur{اگر} \ur{پیادہ} \ur{آخر} \ur{تک} \ur{پہنچا} \ur{تو} \ur{وزیر} \ur{بن} \ur{جائے} \ur{گا۔}
\end{quote}

\begin{quote}
\emph{agar piyādā ākhir tak pahuṅchā to vazīr ban jāe gā.}
\end{quote}

\begin{quote}
If the pawn reaches the end, it will become a queen.
\end{quote}

Masculine.

\begin{cousinweb}
\textbf{\ur{پیادہ}} (\emph{piyādā}) is usually traced to Persian \emph{piyāda}, \textquotedblleft{}on foot; a foot soldier\textquotedblright{}.
\end{cousinweb}

\begin{grammarlens}[title={What you've built}]
One more word for saying what happens if.
\end{grammarlens}

\subsection*{Your turn}
\begin{itemize}
\raggedright
  \item \emph{Say it:} \emph{piyādā}
  \item \emph{Say it:} \emph{piyādā}, once more
  \item \emph{Say it:} \emph{agar piyādā ākhir tak pahuṅchā to vazīr ban jāe gā.} — If the pawn reaches the end, it will become a queen.
  \item \emph{Recall:} say the Urdu for a game board, then the Urdu for a piece (in a board game), then say \emph{piyādā} again
\end{itemize}

\begin{tcolorbox}[breakable,colback=teal!4,colframe=teal!35!black,title={Before you move on}]
What does \ur{پیادہ} mean? (\textquotedblleft{}A pawn (in chess)\textquotedblright{}.) Say it once more.
\end{tcolorbox}
\section[\texorpdfstring{\ur{مات} (māt)}{مات (māt)}]{\ur{مات} (māt) — checkmate}

\label{lesson:UR-C534-mat}



\begin{quote}
Before the new one: say the Urdu for a piece (in a board game), then the Urdu for a pawn.
\end{quote}

\subsection*{You'll want to know: \ur{مات}}
\textbf{\ur{مات}} — \emph{māt} — \textquotedblleft{}checkmate\textquotedblright{}.

\begin{quote}
\ur{اگر} \ur{آپ} \ur{نے} \ur{بادشاہ} \ur{نہیں} \ur{بچایا} \ur{تو} \ur{مات} \ur{ہو} \ur{گی۔}
\end{quote}

\begin{quote}
\emph{agar āp ne bādshāh nahīṅ bachāyā to māt ho gī.}
\end{quote}

\begin{quote}
If you don't save your king, it will be checkmate.
\end{quote}

Feminine.

\begin{cousinweb}
\textbf{\ur{مات}} (\emph{māt}) is usually traced to Persian \emph{māt}, \textquotedblleft{}helpless, beaten\textquotedblright{}, as in \emph{shāh māt}, \textquotedblleft{}the king is helpless\textquotedblright{}.
\end{cousinweb}

\begin{grammarlens}[title={What you've built}]
One more word for saying what happens if.
\end{grammarlens}

\subsection*{Your turn}
\begin{itemize}
\raggedright
  \item \emph{Say it:} \emph{māt}
  \item \emph{Say it:} \emph{māt}, once more
  \item \emph{Say it:} \emph{agar āp ne bādshāh nahīṅ bachāyā to māt ho gī.} — If you don't save your king, it will be checkmate.
  \item \emph{Recall:} say the Urdu for a piece (in a board game), then the Urdu for a pawn, then say \emph{māt} again
\end{itemize}

\begin{tcolorbox}[breakable,colback=teal!4,colframe=teal!35!black,title={Before you move on}]
What does \ur{مات} mean? (\textquotedblleft{}Checkmate\textquotedblright{}.) Say it once more.
\end{tcolorbox}
\section[\texorpdfstring{\ur{لوڈو} (luḍo)}{لوڈو (luḍo)}]{\ur{لوڈو} (luḍo) — ludo (a board game)}

\label{lesson:UR-C534-ludo}



\begin{quote}
Before the new one: say the Urdu for a pawn, then the Urdu for checkmate.
\end{quote}

\subsection*{You'll want to know: \ur{لوڈو}}
\textbf{\ur{لوڈو}} — \emph{luḍo} — \textquotedblleft{}ludo (a board game)\textquotedblright{}.

\begin{quote}
\ur{اگر} \ur{بارش} \ur{ہوئی} \ur{تو} \ur{ہم} \ur{گھر} \ur{میں} \ur{لوڈو} \ur{کھیلیں} \ur{گے۔}
\end{quote}

\begin{quote}
\emph{agar bārish huī to ham ghar meṅ luḍo kheleṅ ge.}
\end{quote}

\begin{quote}
If it rains, we'll play ludo at home.
\end{quote}

Masculine.

\begin{cousinweb}
\textbf{\ur{لوڈو}} (\emph{luḍo}) is borrowed from English \emph{ludo}.
\end{cousinweb}

\begin{grammarlens}[title={What you've built}]
One more word for saying what happens if.
\end{grammarlens}

\subsection*{Your turn}
\begin{itemize}
\raggedright
  \item \emph{Say it:} \emph{luḍo}
  \item \emph{Say it:} \emph{luḍo}, once more
  \item \emph{Say it:} \emph{agar bārish huī to ham ghar meṅ luḍo kheleṅ ge.} — If it rains, we'll play ludo at home.
  \item \emph{Recall:} say the Urdu for a pawn, then the Urdu for checkmate, then say \emph{luḍo} again
\end{itemize}

\begin{tcolorbox}[breakable,colback=teal!4,colframe=teal!35!black,title={Before you move on}]
What does \ur{لوڈو} mean? (\textquotedblleft{}Ludo (a board game)\textquotedblright{}.) Say it once more.
\end{tcolorbox}
